\textbf{Время:} 01:45

\textbf{Тезис:} текущие модельные радиусы и радиусы STAR можно качественно сравнить для одинакового заряда пар, энергии, центральности и LCMS.

\textbf{На экране:} общие условия в шапке; две серии в легенде; панели Rout, Rside, Rlong слева направо; оговорки под графиком.

\textbf{Текст:} Теперь на общих осях стоят мои текущие результаты SMASH и данные STAR. В обоих случаях это пары отрицательных пионов, золото с золотом при трёх ГэВ, центральность 0–10 процентов и радиусы в LCMS.

По легенде находим модельную и экспериментальную серии. Проходим слева направо: в out модельные центры выше экспериментальных; в side последовательности ближе друг к другу; в long модельные центры снова выше. С ростом k-T радиусы в обеих сериях в целом уменьшаются.

Подпись внизу задаёт точность этого сопоставления: центры точек считаны с исходных изображений, а погрешности на общий график не перенесены. Полные отборы тоже пока не согласованы. Поэтому обсуждаем масштабы и зависимости, а точное согласие модели с экспериментом по этому рисунку не устанавливаем.

Следующий шаг для количественного анализа — исходные численные таблицы и согласованные отборы.

\textbf{Переход:} Вернёмся к общей цепочке: что мы узнали, начав с зарегистрированных импульсов?

\textbf{Справка:} 12 текущих модельных центров оцифрованы из красной последовательности kt\_c\_all\_graphs\_neg.jpg. 12 центров STAR для отрицательных пионов взяты без изменения из прежней оцифровки рис. 5 статьи А. Краевой, Universe 10, 188 (2024). Калибровки и исходные пиксели записаны в figures/comparison\_digitized\_current\_2026\_star\_provenance.json; числа — в CSV с тем же префиксом. Бюджет ручного положения маркера не является физической погрешностью.

Общие условия подтверждены автором и статьёй; точные текущие трековые отборы, распределения пар по быстроте, границы всех kT-интервалов и обработка взаимодействий из JPEG не установлены. Нельзя приписывать им настройки 2025 года или вычислять статистическое согласие по центрам без ошибок. Статус STAR Preliminary сохранён. Оцифровка не заменяет исходные численные результаты.

\textbf{Источники:} текущий авторский JPG от 19.10.2026; Kraeva (2024), рис. 5. Старое сравнение SMASH 2025 и STAR сохранено в резерве 47.
